\section{Discussion}
\label{sec:discussion}

\subsection{What the Loop Buys, What the Ladder Prices, and What Neither Can}
The verified loop replaces confidence in the generator with an audited
pipeline, and the evaluation localizes what each stage
contributes. The ladder adds the accounting: which rungs of intelligence
carry measurable weight. The RQ4 results cut both ways -- the
grounding rung is load-bearing, while the top rung's routing benefit
was not resolvable above the deterministic floor -- and the rung's
provider refused assessment precisely on the highest-stakes sessions --
so the LLM came out of that path. A separate rung on the same axis---
the Evidence Judge/filter --- was frozen and tested as a contrast on a
prospective 40-case cohort exercising the actual HTTP and SSH serving
path; on that frozen protocol it produced no detectable difference in
evidence admission or grounded serving relative to the baseline gate
(both arms $0.7500$ unsupported-admission, paired difference $0.0000$,
$p=1.0000$, criterion not met), so it was retained behind its
disabled-by-default boundary as the optional, separately-versioned stage
it is rather than claimed as a contribution. That negative result is
reported in the evaluation and indexed to its frozen artifacts; it is
not averaged into the serving loop that the rest of the paper measures. The loop makes spent intelligence
safe; the ladder keeps it cheap. Neither makes the provider problem
disappear, and the hard-subset census adds a wrinkle: the free-tier
endpoint intermittently refuses security-shaped payloads outright
(HTTP 403 on attacker-tooling states, benign states served). For a
deployment whose critical moments are attack-shaped, provider refusal
is not only a cost risk but a selective availability risk concentrated
exactly where assessment matters most. Every
generative claim in this paper is downstream of a free-tier inference
API: the safety-alignment timeouts that shaped the persona stratum, the
catalog churn the watcher caught, and the quota windows that gated the
study schedule are all the same structural fact. We treat the provider as
an unreliable component and the system degrades safely when it fails, but
a practitioner sizing a deployment should budget for generation being
\emph{eventually unavailable} in exactly the way the fallback path
assumes -- and should read the fold-down rule as the remedy: every
mechanism compiled down to a certifiable equivalent is one fewer
dependency on that provider.

\subsection{Deployment Guidance}
Five lessons from operating the system generalize. First, run the
watcher: provider catalogs drift, and a wired model disappearing silently
is a production incident for any LLM-dependent defense. Second, scope the
HITL gate with care: the gate exists for content classes whose
compromise value is high, and the audit log it produces doubles as the
governance record. Third, keep the deterministic floor maintained: the
fallback transcripts are load-bearing. The system is only as consistent
as its last vetted transcript, and treating those transcripts as second-class
content would hollow out the safety property that makes the LLM path
permissible. Fourth, budget for provider instability as a first-class
operating condition: the watcher's probe history shows the catalog
drifting under observation --- wired models disappeared from the listing
during and after the study, and unwired models appeared in or vanished
from it between consecutive twice-daily probes --- alongside availability
events, a quota-exhaustion window and transient gateway outages. The fallback
path absorbed all of it without serving unverified content, and the
unwired churn never touched the serving path because nothing was wired
to it. An
LLM-dependent defense sized without that budget will be sized wrong.
Fifth, audit the ladder on a schedule: re-price each intelligent rung
against its cheapest certifiable competitor as the deterministic
machinery matures. In this study the audit has already priced two
intelligent rungs and executed one removal -- the kind-aware routing
rung was removed after deterministic fixes reproduced its accuracy
offline, leaving the serving path with one fewer model, provider, and
failure mode, and the retrieval-engine rung is priced against
deterministic BM25 with its removal released as the protocol's next
step rather than claimed (\S\ref{sec:jeveval}) -- and in the released
system the audit is not a manual exercise: the judge-panel audit and
audited stability reseed automatically whenever a new campaign run lands
(Appendix~\ref{app:audit}). A third ladder decision is negative rather
than additive: the Evidence Judge/filter was frozen, tested on a
prospective cohort that exercised the real serving path, and found to
add no measurable value beyond the gate on that configuration, so it
remains an optional, disabled-by-default stage -- a result recorded
against its frozen artifacts rather than folded into the loop that the
evaluation otherwise measures.

One decision we surface rather than bury: the released configuration
ships the timing path \emph{unpadded}, and \S\ref{sec:timing-measured}
shows an attacker can fingerprint the serve path from that channel
(AUC $0.983$) even though modeled padding closes it. The padding
is a deliberate deferral --- a waiting, attacker-visible shell with
homogenized latency changes the attacker's experience of the honeypot
and needs its own deception-effect measurement before it becomes the
default --- but a practitioner whose threat model includes a
latency-fingerprinting adversary should enable it; the measurement
to copy is released.

\subsection{Cost and Practicality}
A defender considering this architecture is entitled to ask what it costs.
The serving host is a commodity box (2~vCPU, ${\sim}4$~GB RAM) shared with
the honeypot listeners; profiling and retrieval are low-cost in-process
steps, and the dominant cost is provider inference, bounded per draft by
output caps and low temperature. The mean served response in the live campaign took
${\sim}44$~s, dominated by the generative turns (the harness derives
the session-duration field from the same per-response statistic; the
campaign's scripted single-turn sessions make them equal by
construction), and the
deterministic strata served in fractions of a second. The economically
interesting number is the fallback rate: when the provider refused or
timed out, the loop degraded to vetted transcripts rather than
open-ended retries, so the cost ceiling per session is two generative
attempts plus one deterministic serve. On a free-tier allowance the full
218-query campaign completes within a daily quota window, which is the
difference between a system a lab can operate and one it cannot.

\subsection{Generalizing the Ladder and the Loop}
Nothing in the loop is specific to honeypots, and neither is the ladder.
The shape, generate under
grounded context, verify the draft against a persisted state structure,
serve only on pass, fall back to vetted content on failure, applies
wherever an LLM writes output whose downstream consumer is adversarial or
safety-critical: SOC copilot drafting incident narratives, retrieval
chatbots answering from a curated corpus, agent pipelines writing into
structured stores. The fiction graph generalizes to any \emph{state
bound} the LLM must respect; the verifier to any checkable draft
property; the deterministic fallback to any canned-but-correct content.
The certifiability axis is the portable part of the ladder: any
mechanism in any such pipeline can be placed on it, priced by marginal
effect, and audited for fold-down as certifiable alternatives mature.
The evaluation here does not establish that the
pattern transfers, but the
architecture section is written so that a practitioner in another domain
can lift the loop wholesale, and the ablation battery (A1--A4 removal
checks, plus the A5 equivalence test) together with the RQ4 pricing
protocol are the template for testing whether each piece they keep
contributes.

\subsection{The Timing Side Channel, Measured}
\label{sec:timing-measured}
The threat model names response timing as the one
channel the observation model leaks; here we measure the leak from the
attacker's seat rather than assert it. Three timing classes are taken
from operational traces: real inter-command gaps from captured sessions
(38 gaps, 13 sessions), LLM-served response latencies from persisted
benchmark rows (376 rows), and the deterministic fast-path timed
in-process on the production template route (384 calls, median $0.1$~ms).
A real-vs-real drift-floor control (single-gap directional AUC $0.676$)
sets the sensitivity baseline. Three findings, all against that floor.
The unmitigated deployment -- a mix of mostly LLM serves with a small
deterministic fraction -- is fingerprintable at AUC $0.983$: a passive
attacker can classify the serve path from timing alone. An all-deterministic
surface is \emph{equally} detectable (AUC $1.0$, the ``too fast''
extreme): maximizing fast-path coverage by itself is not a mitigation
but a different fingerprint. Padding the fast-path's response times
with additive sleep drawn from the real inter-command-gap envelope,
modeled analytically rather than measured as deployed padding, closes
the channel: AUC $0.551$, at or below the drift floor. The deployment
guidance is therefore two-part and quantified: raise deterministic
coverage \emph{and} shape its timing to that envelope; LLM latency
cannot be padded down, so the
mitigation direction is fixed by physics, not preference. The serving
path that carries these numbers is the released default: a commodity
host (2~vCPU, ${\sim}4$~GB RAM), the deterministic floor maintained as a
first-class content class, and the generator treated as an unreliable
component whose failures degrade to vetted transcripts rather than
open-ended retries. The same stochasticity that makes provider inference
a cost also makes it an availability risk; the loop is built so that the
cost is bounded per draft and the availability risk is absorbed by a
path that is correct by construction rather than by hope. The
host-fiction layer (\S\ref{sec:hostfiction}) is timing-neutral by the same
paired measurement: byte-identical serves with the store on and off.
The system deploys deception against real attackers, which obligates
care. All credential material shown to attackers is registered canary
secrets: nothing real is ever exposed, and any canary appearing in
outbound traffic is attributable to bait interaction by construction. The
HITL gate keeps a human decision between the system and the highest-stakes
content classes, and the append-only audit log records every gate
decision. The testbed is isolated from production networks, and the
deception targets only actors actively attacking it. We release scenario
sets and scoring harnesses so that the community can audit the
evaluation, not just the code.

\subsection{Limitations}
The deployment observes tailnet traffic only; the ambient-internet axis
(scanners, botnets, opportunistic ransomware) is absent by design choice
and limits ecological claims. The consensus study's three evaluator
personas (models, not human raters) share one judge model family; multi-model judging failed on the
provider at study time (a listed model returned 404 on real calls despite
catalog presence), a gap we flag for the watcher's future checks. The
$0.750$ agreement figure should therefore be read as within-model
consistency -- how stably one model's prompted personas concur -- rather
than inter-annotator agreement across independent judges. The shell-scope listener constrains web-shaped deception scenarios, as the
weakest consensus scenario illustrates. And the persona-realism stratum,
while stable run-to-run, remains the stratum where judge variance lives;
its scores should be read with the MAD in \S\ref{sec:stability} in mind.
The Evidence Judge/filter and the conditional Laya scorer are separate
experiment identities --- a frozen deterministic stage tested on the
prospective cohort and a separately-versioned feasibility study that did
not pass its gate --- and neither is claimed as part of the served loop's
measured contributions; their records stand where they stand.
The RQ4 evidence carries its own weight limits: the real-log assessment
sample is 13 sessions, the closed-loop replay pairs number five
clean families drawn from seven attempted (two excluded for attacker-LLM
provider errors), and the simulated attacker, however instrumented, is
not an adversary with a stake in the outcome. The routing fold-down rests on
200 held-out sessions that sit at ceiling with no assessment output
available on any of them (\texttt{jev\_available}$=0$ in the released
artifact), plus a full-corpus census locating the
informative stratum (six sessions, one rescued live, five lost to
provider refusals); it is stated at the benchmark's resolution rather
than as general subsumption. The deception-effect claims for the
top rung (bait exposure, the dwell null) are directional, and we price
them as such rather than as established effects.
